\section{The End}

This paper positioned a discrete, deterministic, reversible base against the continuum programs in one
disciplined frame, and showed, one result at a time, where the base derives a continuum result, where it runs parallel to a
program by a different route, and where the bridge is a well-posed open conjecture with a stated falsifying test. The firm
floor is exact and integer-checked. Charge conservation is the continuity law $\nabla\!\cdot J = 0$, shown on a flat torus
and on the genuinely curved $\{3,4,3,4\}$ mesh. The 24-cell coin is selected by a spinor test and carries the isotropy a
four-dimensional lattice gas needs. One recoverability functional carries horizons, decoherence, and conservation
together. The middle rungs run parallel, and the open walls are listed with their tests.

The conjecture is this. The coarse-grained face of one growing mesh of one substance, carrying a ternary value, moved by
one local reversible charge-conserving rule, is a conserved-current field theory of the kind the continuum programs
describe. That claim is not proven. It is supported by a firm anchor and a clear list of unbuilt steps, and it stands or
falls on experiments hitting thresholds fixed before they ran.

This is exploratory, feasibility-stage work, built to be tested and falsified. The whole comparison rests on a reproducible suite of deterministic, controlled experiments, each paired
with a control that would break it if the result were an artifact \cite{cluesurf2026repo}. The invitation is to run them,
to push the parallel rungs toward measured ones, and to test the open walls against nature. A theory should derive what it
can, state plainly what it cannot, and remain open to being shown wrong.
